\chapter{Conclusion}
\label{ch:conclusion}

This thesis asked whether the IVM$^+$ approach of Abo~Khamis et al.\ is beneficial in practice when realised at the SQL level on an existing incremental engine. We built a query decomposer that rewrites a join query into the bag views of a tree decomposition, in both IVM$^+$ and an adjusted variant, with a semiring extension for aggregates, and evaluated the rewrites on Feldera against a flat-join baseline across graph, TPC-H and IMDB workloads.

The central finding is that the decomposition helps exactly where the flat query blows up: on cyclic queries over dense data. On \texttt{cit-HepTh} the baseline maintains Cyclic Query~2 in $80.77$\,min, while the decomposition maintains the same $266$-million-row result in under two minutes, a $42.96$-fold speedup, and at essentially the baseline's own peak memory of about $16$\,GiB. Our running example, the diamond $Q_\diamond$ (Cyclic Query~3), has its full reconstruction maintained $5.6$ to $6.7$ times faster on the two dense graphs, or $9.5$ to $12.7$ times faster without the reconstruction. On \texttt{cit-HepTh} Cyclic Queries~1, 4 and 5 the baseline does not finish within three hours at all, whereas the decomposition does. Notably, these gains hold without Feldera implementing the worst-case-optimal joins that the $O(N^{w(Q)})$ bound of Section~\ref{sec:fhw} assumes: the advantage comes from the bag structure constraining the query, not from that bound.

The benefit is not universal. Where the flat query does not build large intermediates, the decomposition only adds overhead. On the acyclic TPC-H 4 and 5 Table Path Queries, whose foreign-key to primary-key joins leave nothing for the bags to filter, it matches the baseline's throughput at best and uses up to $2.6$ times its memory. On the sparse \texttt{p2p-Gnutella31} graph and on TPC-H Cyclic Query~7, where the baseline stays small, the decomposition is likewise slower. The advantage is specific to cyclic queries over dense data, not to decomposition in general.

Among the variants, the reconstruction is the expensive step. Stopping at the bag hierarchy rather than rejoining it into the full result is markedly cheaper, on \texttt{cit-HepTh} Cyclic Query~2 $0.44$ against $1.88$\,min, and the semiring extension recovers the result's size at nearly the same low cost, returning counts as large as $234$ billion rows, at under $13$\,GiB in that case, where materialising the full result is infeasible. Between the two decomposition forms, our variant, which drops the projected-atom semijoins, is the better default, no slower than IVM$^+$ on every query but TPC-H Cyclic Query~7 and consistently lighter. The tree's orientation also matters, by up to nearly a factor of four and sometimes deciding whether a query finishes, but which orientation is best depends on the data and is not something our measurements predict.

These results are specific to Feldera and its DBSP engine, and a different engine, in particular one with worst-case-optimal joins, might shift the boundaries. Correctness was checked by comparing the output count across all variants of each query, a strong but not exhaustive guarantee. Future work could extend the semiring aggregation beyond \texttt{COUNT} to further aggregates, measure the intermediate bag sizes directly to explain the orientation results our current metrics leave open, and widen the decomposer to the full query surface. Even so, the core question has a clear answer: realised as plain SQL views on an unmodified engine, the IVM$^+$ decomposition is a worthwhile rewrite for cyclic queries over dense data, the queries where the flat plan would otherwise blow up.
